\documentclass[10pt,a4paper]{article}
\usepackage[T1]{fontenc}
\usepackage[margin=23mm]{geometry}
\usepackage{lmodern,amsmath,amssymb,amsthm,mathtools,microtype}
\usepackage[hidelinks]{hyperref}
\newtheorem{theorem}{Theorem}
\newtheorem{lemma}[theorem]{Lemma}
\title{An actual-factor socle certificate for Erd\H{o}s--Straus}
\author{The Clankers}
\date{15 September 2026 --- research preprint}
\begin{document}\maketitle
\begin{abstract}
A characteristic-two group algebra can force every unit residue to occur among
an explicitly bounded family of actual divisors. We give a complete binary-prefix
test for one leading-socle selection class, preserve original factor multiplicities
through a positive binary-submask lift, and return an ordered ES solution.
The construction allows new primes to replace missing powers of three in the
preceding seed method. A hard-prime example has exactly one three-factor/cofactor
state. The group-algebra strategy is classical; no universal ES or priority claim
is made.
\end{abstract}

\paragraph{Original arithmetic and the allowed modulus reduction.}
Let $k\ge4$, $u=2^{2k-5}$, $m=2^k$, $p\equiv1\pmod8$, $p>2u$, and
$N=p+4u=\prod_q q^{e_q}$, with actual prime exponents. Put
$h=k-1$ when $p\equiv1+2^{k-1}\pmod m$, and $h=k$ otherwise; set
$n=2^h$, $L=2^{h-2}$ and $r_0=h-2$.
A divisor $d\mid N$ in class $-1\pmod n$ gives $Q=d$, except that in the
half-turn case its full residue may be $-p\pmod m$: then put $R=d$ and
$Q=N/d$. In all cases $RQ=N$, $Q\equiv-1\pmod m$.
In the half-turn case complementation exchanges the two distinct lift classes;
the map has two preimages per original state. Otherwise it is bijective.

The inherited middle-state return is
\[
 a=(p+R)/4,\quad d_0=\gcd(a,u),\quad
 (h_0,r,s)=(d_0^2/u,u/d_0,a/d_0),\quad \lambda=(r+s)/R,
\]
\[
 (x,y,z)=(a,ph_0s\lambda,ph_0r\lambda),\qquad
 u=pa^2/(Ry-pa).                                             \tag{1}
\]
Indeed, $u\mid a^2$ is equivalent to $2^{k-2}\mid a$, obtained from the
cofactor congruence. Both factors are $3\pmod4$, at least three and less than
$p$, since $N<3p$. Prime valuations give $a=h_0rs$, $u=h_0r^2$ and
$\gcd(r,s)=1$. The middle gate gives $R\lambda=r+s$, and clearing the
reciprocal sum proves $4/p=1/x+1/y+1/z$. These Type II coordinates and the
fixed-seed atlas are inherited \cite{ET,Prior}.

\paragraph{Actual divisor words before parity reduction.}
For a selected budget $0\le f_q\le e_q$, keep
\[
 B_f=\prod_q\prod_{a:\,\mathrm{bit}_a(f_q)=1}(1+[q]^{2^a})\in\mathbb Z[G],
 \qquad G=(\mathbb Z/n\mathbb Z)^\times.                      \tag{2}
\]
Every term is an original divisor: its exponent at $q$ is a binary submask
of $f_q$. Unique factorization and the binary expansion invert the term map.
Coefficients after residue collision retain the number of such terms.
Modulo two, Frobenius gives $\overline B_f=\prod_q(1+[q])^{f_q}$.
A nonzero parity proves positive occupancy. Zero parity may be a supported
even coefficient, not absence; (2) is not erased by reduction.

For $g=3$ or $5$, $G=\langle-1\rangle\times\langle g\rangle$ with
$|\langle g\rangle|=L$. Induction gives the orders modulo $2^h$; the two
subgroups are respectively $1,3\pmod8$ and $1,5\pmod8$. The substitutions
\[
 S=1+[-1],\quad T=1+[g],\qquad
 [(-1)^\epsilon g^j]\longleftrightarrow(1+S)^\epsilon(1+T)^j
\]
give an algebra isomorphism $\mathbb F_2[G]\simeq\mathbb F_2[S,T]/(S^2,T^L)$,
with the displayed inverse. Its joint annihilator is the line
\[
 \operatorname{Ann}(S,T)=\mathbb F_2ST^{L-1},\qquad
 ST^{L-1}=\sum_{t\in G}[t].                                \tag{3}
\]
The monomial basis proves the annihilator statement; the odd binomial
coefficients of $(1+[g])^{L-1}$ prove the last identity. Highest-augmentation
products forcing occupied coefficients are a classical group-algebra method,
not a new principle here \cite[\S1, pp.~2--3]{Petrov}.

\newpage
\begin{theorem}[Actual-factor selection]
Write $q=(-1)^{\epsilon_q}g^{j_q}\pmod n$, $0\le j_q<L$, and set
$w_q=2^{v_2(j_q)}$ for $j_q\ne0$, with sentinel $w_q=L$ at $j_q=0$.
Suppose an actual selection $f_q\le e_q$ and a pivot $m_0=2^a$ satisfy
\[
 \sum_qw_qf_q=L-1+m_0,\quad
 \sum_{\epsilon_q=1,w_q=m_0}f_q\equiv1\pmod2,\quad
 \epsilon_q=1,
w_q>m_0\Longrightarrow f_q\text{ even}.       \tag{4}
\]
Then $\overline B_f=ST^{L-1}$, every unit residue has an original divisor
representative in (2), and (1) returns a positive integral ES witness.
At most $L$ prime occurrences are needed in one marked factor $R$ or $Q$.
\end{theorem}
\begin{proof}
For $j_q\ne0$, write
\[
 1+[q]=A_q+\epsilon_qSB_q,\quad
 A_q=1+(1+T)^{j_q}=T^{w_q}U_q,\quad U_q(0)=B_q(0)=1.
\]
The part without $S$ vanishes by its degree. The coefficient of $S$ is
\[
 \sum_{\epsilon_q=1,\ f_q\text{ odd}}
 B_qA_q^{f_q-1}\prod_{t\ne q}A_t^{f_t}.
\]
Its summand has degree $L-1+m_0-w_q$. Larger weights have even selected
exponent; smaller weights vanish modulo $T^L$; pivot summands each equal
$T^{L-1}$ and occur in odd number. A selected identity-residue prime is
excluded by the weight equation. A selected $j_q=0$ outside prime forces
$m_0=L$ and exactly one such occurrence; its factor is $S$ and the remaining
weight is $L-1$, proving that exceptional case directly. Removing one pivot
occurrence leaves weight $L-1$, so there are at most $L$ occurrences total.
Equations (2)--(3) and the cofactor return prove the claim.
\end{proof}

The weights can be computed without a discrete-log search. For
$x=(-1)^{\epsilon_q}q$, use $w_q=1$ when $x\equiv g\pmod8$,
$w_q=2^{v_2(x-1)-2}$ when $x\equiv1\pmod8$ but not modulo $n$, and $w_q=L$
when $x\equiv1\pmod n$. The order identities prove these formulas.

\begin{theorem}[Exact finite capacity test for (4)]
For each weight $2^b$ let
\[
 I_b=\sum_{\epsilon_q=0,w_q=2^b}e_q,\quad
 O_b=\sum_{\epsilon_q=1,w_q=2^b}e_q,\quad
 P_b=\sum_{\epsilon_q=1,w_q=2^b}\lfloor e_q/2\rfloor.
\]
For a pivot $a$ with $O_a>0$, put, for $0\le b<r_0$,
\[
 c_b=I_b+\mathbf1_{b<a}O_b
 +\mathbf1_{b=a+1}\lfloor(O_a-1)/2\rfloor
 +\mathbf1_{b\ge a+2}P_{b-1}.                               \tag{5}
\]
A selection (4) exists exactly when
\[
 \boxed{\sum_{b=0}^{j-1}2^bc_b\ge2^j-1\quad(1\le j\le r_0).} \tag{6}
\]
Thus checking all pivots gives a terminating arithmetic sufficient test for
arbitrary original prime-factor multiplicities.
\end{theorem}
\begin{proof}
Reserve one pivot occurrence. The remaining odd pivot aggregate is grouped
into pairs of weight $2^{a+1}$; higher outside primes must pair their own
occurrences, giving $P_b$, not $\lfloor O_b/2\rfloor$. Lower outside and all
inside occurrences are unrestricted. The remaining target weight is $L-1$,
with exactly the coin capacities (5). Necessity of (6) follows modulo $2^j$.
For sufficiency, at each new denomination $2^b$, lower coins cover a complete
interval at least through $2^b-1$. Its translates by all available multiples
of $2^b$ touch or overlap, so subset sums fill every integer through the new
total. The last inequality includes $L-1$. Descending greedy allocation and
bounded compositions among the original primes give a full inverse selection.
\end{proof}
The test is exact for the leading-term class (4), not necessary for every ES
witness or every possible group-ring cancellation.

A quantitative consequence retains unused factor capacity. Put
$M_f=\prod_q\lfloor(e_q+1)/(f_q+1)\rfloor$. The translates of (2) by
$\prod_q q^{(f_q+1)t_q}$, with $0\le t_q<\lfloor(e_q+1)/(f_q+1)\rfloor$,
are disjoint by quotient and remainder. Each contains a coarse target.
Therefore the original state count is at least $M_f$ when $h=k$, and at
least $\lceil M_f/2\rceil$ in the half-turn case.

\newpage
\paragraph{A mixed contribution with a unique arithmetic return.}
At the deterministically checked hard prime
\[
 p=1794769\equiv529\pmod{840},\quad k=5,\quad u=32,\quad
 N=3^2\cdot13\cdot23^2\cdot29,
\]
we have $h=4$, $L=4$. In the $g=3$ basis, the pairs $(\epsilon,w)$ for
$3,13,23$ are $(0,1),(1,1),(1,2)$. The actual selection
$(f_3,f_{13},f_{23})=(2,1,1)$ satisfies (4) at pivot two. Its positive lift is
\[
 (1+[9])(1+[13])(1+[23]),
\]
whose eight integer residue coefficients modulo 16 are all one.
The divisor $207=9\cdot23$ is $-1$ modulo 16, but $-p$, not $-1$, modulo
32. Thus it is the residual; $Q=8671=13\cdot23\cdot29$.
The original marking is
\[
 (a,u,h_0,r,s,\lambda)=(448744,32,2,4,56093,271),
\]
\[
 \boxed{\frac4{1794769}=\frac1{448744}
 +\frac1{54565295814214}+\frac1{3891059192}.}
\]
All 36 divisors of $N$ give exactly this one original state. Neither target
has a one- or two-prime word. The preceding $3$-power threshold required
three occurrences but only two are present. The two simpler tests using
inside factors in either $\langle3\rangle$ or $\langle5\rangle$ alone also
fail. The mixed contribution supplies the missing arithmetic capacity.

For $p=165601$, $k=4$, $N=3\cdot13\cdot31\cdot137$, the weights of the
inside primes $3$ and $137$ are one and two. The outside prime $13$ completes
the socle selection. Its return is
$(41408,1145153436736,221242936)$, and the disjoint-packing bound gives
both original states. The prime $31$ already supplies an alternative cofactor;
this example illustrates the weighted mechanism, not failure of every simpler
construction. At $p=852769$, $k=6$, prime $17$ contributes weight four beside
three actual powers of $3$; outside prime $13$ gives the unique state
$(213280,863923218520,518483552)$. It requires no sevenfold $3$-power block.

\paragraph{Scope, finite evidence, and the remaining problem.}
The complete test through two million contains 4519 hard-class primes,
37289 valid dyadic branches and 4583 original states in 3097 occupied branches.
The previous sufficient capacity theorem covers 964 primes. The new prefix
test covers 1194 primes; their union covers 1196. This adds 232 primes to
that test, not to known overall ES coverage. Two previously certified primes
remain outside this new sufficient test and are not discarded.

Every tested positive certificate is returned to actual factors, the full
modulus orientation, the original $a$ factorization, and the ordered
reciprocal identity. The main verifier uses valuations and prefix inequalities;
a separate implementation uses finite logarithms and parity-aware bounded
knapsack and agrees on the complete branch digest. The source, exact JSON,
full proofs and typed maps accompany this note. No Lean or independent
mathematical review is claimed.

The final universal step is still arithmetic. The construction proves
\emph{when actual factors form a nonzero top product}; it does not assert
that every prime has a seed with that capacity. If every pivot fails, a
specific finite binary-prefix deficit remains. Cross-seed constraints must
force one such deficit to disappear or produce a different positive selection.
A supported zero coefficient, a generated subgroup, and an available bounded
factor word remain distinct observations throughout.

\begingroup\small
\begin{thebibliography}{3}
\bibitem{ET} Elsholtz, C., \& Tao, T. (2013). Counting the number of solutions
to the Erd\H{o}s--Straus equation on unit fractions.
\emph{Journal of the Australian Mathematical Society, 94}(1), 50--105.
\S2, Propositions 2.6--2.7; arXiv:1107.1010v6.
\bibitem{Petrov} Petrov, F. (2020). \emph{Combinatorial results implied by many
zero divisors in a group ring}. arXiv:1606.03256v4. \S1, pp.~2--3.
The general augmentation method and its Olson antecedent are attributed there;
no first-use claim for that method is made here.
\bibitem{Prior} The Clankers. (2026). \emph{All-phase cofactor bounds, valuation
correlations and logarithmic escape}. 15 September 2026 source tranche,
Sections 1--3; predecessor fixed-seed PR8, commit \texttt{79054f8b\ldots{}ba750}.
Exact revisions and read scopes are in the accompanying source ledger.
\end{thebibliography}\endgroup
\end{document}
